% 摘要
\chapter*{摘\quad 要}
\addcontentsline{toc}{chapter}{摘要}

高比例可再生能源并网背景下，净负荷（总负荷减去风光出力）的准确预测与物理可解释性对电力系统调度至关重要。Kolmogorov-Arnold网络（KAN）通过边上可学习的B样条函数实现"先学习、后提取"的符号回归路径，有望同时满足预测精度与公式可解释性的需求。然而，将KAN符号提取应用于包含自回归滞后特征的时间序列预测任务时，会出现系统性的"公式坍缩"现象：稀疏化-剪枝过程将全部物理气象变量（风速、辐照度、温度等）剪除，最终提取的公式退化为仅含负荷滞后项的极简表达。

本文以ARPA-E PERFORM数据集的5分钟采样净负荷数据为对象，围绕上述公式坍缩现象的成因与应对，开展了四项递进实验。首先，构建KAN教师-符号提取管线并验证其预测有效性；在当前标准对齐的主实验设置下，KAN教师的技能分数达到0.453，显著优于MLP基线（0.430，$p=0.0005$）和PySR（0.076）。其次，提出变量进入率（VER）和公式出现率（FAR）等可辨识性诊断指标，系统记录了直接净负荷目标下9组符号配置全部坍缩（VER=0/9）的现象。第三，设计S2自回归阻断干预实验，以干预对照而非仅凭观察性相关来检验"自回归捷径竞争"假说：阻断滞后特征后，聚焦风电任务的$\Delta\VER=+0.60$（95\% CI=[0.20, 1.00]）；而在严格matched的直接净负荷任务上，4个有效paired seed仅给出$\Delta\VER=+0.25$（95\% CI=[0.00, 0.75]），另有1个seed两次未产出post-prune checkpoint。这表明shortcut competition的干预性证据在focused wind设置下最强，在direct net-load设置下则更弱且不稳定。第四，设计S3结构化分解方案，将净负荷拆分为负荷、风电、光伏三个子任务。S3结构化组合预测器在保持预测精度并进一步提升组合表现（skill=0.515）的同时实现了全部子任务FAR=3/3，表明结构化分解后可以稳定提取具有物理含义的局部公式；进一步地，本文已经完成净负荷总公式的显式合成与test复算，但复算结果的RMSE=2644.36、skill=-0.021，说明总公式对象相对结构化组合预测器仍存在明显transfer gap。

本文支持如下判断：在当前实验与数据处理设置下，自回归捷径竞争比“变量无关性”更能解释KAN符号提取中物理气象变量被剪除的现象。S2阻断实验为这一解释提供了干预性支持，S3分解方案则给出了兼顾预测精度与局部公式可辨识性的建设性路径。研究同时诚实报告了方法的边界条件：S2阻断在直接净负荷任务上带来明显精度代价，太阳辐照在长期时域外推退化，风电可辨识性随时域呈非单调变化；尽管S3总公式已经完成显式合成与test复算，但其相对结构化组合预测器仍存在明显transfer gap，因此预测器对象与公式对象必须分开表述。

\vspace{1em}
\noindent\textbf{关键词：}KAN；符号回归；净负荷预测；可解释性；自回归竞争；可辨识性

\newpage

% Abstract
\chapter*{Abstract}
\addcontentsline{toc}{chapter}{Abstract}

Accurate and physically interpretable forecasting of net load (total demand minus wind and solar generation) is essential for power system dispatch under high renewable penetration. Kolmogorov-Arnold Networks (KAN) offer a ``learn-then-extract'' symbolic regression pathway through learnable B-spline activation functions on network edges, potentially satisfying both predictive accuracy and formula-level interpretability. However, when applied to time-series forecasting tasks containing autoregressive lag features, KAN symbolic extraction exhibits systematic ``formula collapse'': the sparsification-pruning pipeline eliminates all physical meteorological variables (wind speed, irradiance, temperature, etc.), and the extracted formula degenerates into a minimal expression containing only load lags.

Using 5-minute net load data from the ARPA-E PERFORM dataset, this thesis conducts four progressive experiments to investigate the cause and remedy of formula collapse. First, a KAN teacher-to-symbol extraction pipeline is established and validated; within the current canonical matched main experimental setting, the KAN teacher achieves a skill score of 0.453, significantly outperforming an MLP baseline (0.430, $p=0.0005$) and PySR (0.076). Second, identifiability diagnostics---Variable Entry Rate (VER) and Formula Appearance Rate (FAR)---are proposed, documenting complete collapse (VER = 0/9) across all nine symbolic configurations on the direct net load target. Third, an S2 autoregressive blocking intervention compares blocked and unblocked settings and provides interventional support for the shortcut-competition hypothesis: after blocking lag features, $\Delta$VER = +0.60 (95\% CI = [0.20, 1.00]) for the focused wind task, while the strict matched direct net-load setting yields only $\Delta$VER = +0.25 (95\% CI = [0.00, 0.75]) across four valid paired seeds, with one additional seed twice failing to produce a post-prune checkpoint. This indicates that the interventional evidence for shortcut competition is strongest in the focused wind setting and weaker, less stable in the direct net-load setting. Fourth, an S3 structured decomposition scheme partitions net load into load, wind, and solar sub-tasks. The structured composite predictor preserves predictive accuracy while improving the final combined result (skill = 0.515) and achieving FAR = 3/3 for all sub-tasks, and the explicit replayed net-load formula is now closed on the test set. However, the replayed total formula reaches RMSE = 2644.36 and skill = -0.021, revealing a clear transfer gap relative to the structured predictor.

Overall, the evidence supports the following interpretation: under the current experimental and preprocessing setup, autoregressive shortcut competition---rather than variable irrelevance---more plausibly explains why physical meteorological variables are pruned during KAN symbolic extraction. The S2 blocking experiment provides interventional support for this interpretation, and the S3 decomposition offers a constructive solution balancing accuracy and interpretability. Boundary conditions are honestly reported: direct blocking on the net-load task carries a noticeable accuracy cost, solar identifiability degrades at long-horizon extrapolation, wind identifiability varies non-monotonically across forecast horizons, and although the S3 total formula has been explicitly closed and replayed, it still exhibits a substantial transfer gap relative to the structured predictor.

\vspace{1em}
\noindent\textbf{Keywords:} KAN; symbolic regression; net load forecasting; interpretability; autoregressive competition; identifiability
